\subsection{Characterizations of strict morphisms}

PROPOSITION 1. --- Let E and F be two locally convex spaces and u a continuous linear mapping from E into F. In order that u be a strict morphism, it is necessary and sufficient that the following condition be satisfied :

(MS) For every semi-norm \(p \in \mathrm{S}(\mathrm{E})\) , which is null on the kernel of \(u\) , there exists \(q\) in \(\mathrm{S}(\mathrm{F})\) such that \(p \leqslant q \circ u\) .

Let N be the kernel and M the image of u; we introduce the canonical decomposition of u, let

\[
\mathrm{E} \xrightarrow {\pi} \mathrm{E} / \mathrm{N} \xrightarrow {\tilde {u}} \mathrm{M} \xrightarrow {\iota} \mathrm{F}.
\]

The continuous semi-norms on E which are null on N, are the semi-norms \(p_{1} \circ \pi\) where \(p_{1}\) ranges over S(E/N); similarly S(M) consists of the semi-norms \(q_{1}\) for which there exists \(q \in \mathrm{S}(F)\) with \(q_{1} \leqslant q/F\) . Finally, u is a strict morphism if and only if the bijective continuous linear mapping \(\tilde{u}\) has a continuous inverse; this also means that every semi-norm in S(E/N) is of the form \(q_{1} \circ \tilde{u}\) with \(q_{1}\) in S(M). Prop. 1 follows immediately from these remarks.

PROPOSITION 2. --- Let E and F be two Hausdorff locally convex spaces and u a continuous linear mapping from E into F. In order that u be a strict morphism, it is necessary and sufficient that its transpose \({}^t u\): F' → E' satisfy the following conditions :

\begin{enumerate}
\def\labelenumi{\alph{enumi})}
\item
  The image of \(^{t}u\) is closed in \(E'\) for \(\sigma(E', E)\) .
\item
  Every equicontinuous subset of \(\mathbf{E}'\) , contained in the image of \({}^t u\) is the image under \({}^t u\) of an equicontinuous subset of \(\mathbf{F}'\) .
\end{enumerate}

If this is so, we have \(\operatorname{Ker}^t u = (\operatorname{Im} u)^\circ\) and \(\operatorname{Im}^t u = (\operatorname{Ker} u)^\circ\) and there exist canonical isomorphisms from \(\operatorname{Coker}^t u\) onto the dual of \(\operatorname{Ker} u\) and from \(\operatorname{Ker}^t u\) onto the dual of \(\operatorname{Coker} u\) .

Let N be the kernel and I the image of u. By cor. 2 of II, p.~47, the kernel of \(^{t}u\) is the orthogonal of I, and the closure of the image of \(^{t}u\) for \(\sigma(E', E)\) is the orthogonal \(N^{\circ}\) of N. The conjunction of a) and b) is then equivalent to the following condition:

b') Every equicontinuous subset of \(E'\) contained in \(N^{\circ}\) is the image under \(^{t}u\) of an equicontinuous subset of \(F'\) .

Since \(N^{\circ}\) can be identified with the dual of E/N, prop. 9, (i) of IV, p.~8, shows that the equicontinuous subsets of \(E'\) contained in \(N^{\circ}\) are the sets which are contained in a set of the form \(H_{p}\) , where p is a continuous semi-norm on E, vanishing on N. The condition \(b')\) then says that, for every semi-norm \(p \in \mathrm{S}(\mathrm{E})\) which is null on N, there exists \(q \in \mathrm{S}(\mathrm{F})\) such that \(\mathrm{H}_{p} \subset {}^{t}u(\mathrm{H}_{q})\) . By Hahn-Banach theorem (II, p.~23, cor. 1 and 2, p.~63, th. 1 and cor. 1), we have \(u(\mathrm{H}_{q}) = \mathrm{H}_{q \circ u}\) , and the relations \(H_{p} \subset H_{p'}\) and \(p \leqslant p'\) are equivalent for all semi-norms p and \(p'\) in S(E). Consequently, the relation \(H_{p} \subset {}^{t}u(\mathrm{H}_{q})\) is equivalent to the relation \(p \leqslant q \circ u\) . By prop. 1, property \(b')\) implies that u is a strict morphism.

Suppose that \(u\) is a strict morphism. We have already seen that the kernel of \(^t u\) is the orthogonal of I and the image of \(^t u\) is the orthogonal of N. The cokernel of u is the space F/I and its dual can be identified with \(I^{\circ} = Ker^{t}u\) . Similarly, the dual of N = Ker u can be identified with \(E^{\prime}/N^{\circ}\) (IV, p.~8), i.e.~with the cokernel of \(^{t}u\) since \(N^{\circ}\) is the image of \(^{t}u\) .

Remark. --- With the notations of prop. 2, property \(b'\) ) also implies that \(u\) is a strict morphism for the weakened topologies (II, p.~49, cor. 3).

PROPOSITION 3. --- Let E and F be two locally convex spaces, and u a continuous linear mapping from E into F. We assume that E is Hausdorff and that F is metrizable. For u to be a strict morphism, it is necessary and sufficient that the image of \({}^t u\) be closed in E' for the weak topology \(\sigma(E', E)\) .

The necessity follows from prop. 2.

Conversely, suppose that the image of \(^t u\) is closed for \(\sigma(E', E)\) and introduce the canonical decomposition of \(u\) as in the proof of prop. 1. By the above remarks, the inverse mapping \(\tilde{u}^{-1}\) of \(\tilde{u}\) is continuous for the weakened topologies. But the subspace \(M = u(E)\) of F is metrizable, hence bornological (III, p.~12, prop. 2); consequently (IV, p.~7, prop. 7, (ii)), \(\tilde{u}^{-1}\) is continuous, hence \(u\) is a strict morphism.
% semantic-fragments: legacy-input-seam
\relax{}
